% Chapter 2: Literature Survey (Massive Expansion for 80+ Page Project)
\chapter{LITERATURE SURVEY}
%% \thispagestyle{plain} removed

Over the last 2 decades, there has been a significant growth in the cloud management industry. In order to gain insight into where Nebula fits into the existing marketplace, we carried out a comprehensive survey of three different types of tools Public Cloud Native Orchestrator, Third-Party IaC Tools, and Emerging Unified Platforms.

\section{Critical Analysis of 10+ Cloud Management Tools}

\subsection{Cloud-Native Orchestrators (Provider-Specific)}
1. \textbf{AWS CloudFormation}: CloudFormation is a powerful, JSON/YAML based service which enables declarative definition of infrastructure in AWS. The key benefit of it is that it is highly integrated with AWS IAM and introduction of new services. But the biggest disadvantage is its walled garden characteristic, as it cannot deal with resources in Azure or GCP. \\
2. \textbf{Azure Resource Manager (ARM)}: Azure templates are used to define infrastructure using ARM templates. Though it provides very powerful "Blueprints" to compliance, its syntax is also notoriously verbose and difficult to debug, when compared to Terraform. \\
3. \textbf{GCP Deployment Manager}: Google Cloud native IaC service. It supports Python and Jinja2 templates, providing more programmability than CloudFormation but with no huge support of third-party tools.

\subsection{Cross-Cloud Infrastructure as Code (IaC) Tools}
1. \textbf{Terraform (HashiCorp)}: Terraform is the standard in the industry and the heart of Nebula engine. Its Provider architecture enables it to communicate to any API. Non-DevOps engineers, however, tend to be scared away by its CLI-only status and complexity of State Management. \\
2. \textbf{Pulumi}: A relatively new tool, a competitor to Terraform, which uses standard pro- gramming languages (Python, Go, Node.js). Although it is flexible, its ecosystem of predefined modules is small in comparison to Terraform. \\
3. \textbf{Ansible}: Is mainly a configuration management tool, but often used for provisioning. It is less appropriate to the Infrastructure Lifecycle management (e.g., the time when a resource has been deleted or changed) because of its imperative nature.
 \\
4. \textbf{Chef \& Puppet}: Legacy configuration management tools that have been mostly replaced by containerization (Docker) and IaC (Terraform).


\subsection{Unified Multi-Cloud Management Platforms}
1. \textbf{VMware Cloud (VMC)}: Offers a uniform vSphere experience on other clouds. Its main disadvantage is that it is very expensive to license as an enterprise, and is therefore unavailable to small developers.
2. \textbf{Google Anthos}: This is a managed service with which it is possible to manage Kubernetes in both AWS and Azure. It is however, purely containerized workload-oriented and there is no general-purpose AWS or Azure VM/Storage management.
3. \textbf{Morpheus Data}: A commercial "Cloud Management Platform" (CMP). Although potent, it is a Black Box that can be hardly extended to custom AI-based troubleshooting, and that is where Nebula offers its own value.

\section{Summary of Comparative Analysis}

\begin{table}[H]
\centering
\caption{Comparison of Cloud Management Tools with Nebula}
\label{tab:comparison}
\begin{tabular}{|l|l|l|l|l|}
\hline
\textbf{Feature} & \textbf{Native Tools} & \textbf{Terraform} & \textbf{Anthos} & \textbf{Nebula} \\ \hline
Multi-Cloud Support & No & Yes & Partial & \textbf{Full} \\ \hline
Unified GUI & No & No & Yes & \textbf{Yes} \\ \hline
AI Error Remediation & No & No & No & \textbf{Yes} \\ \hline
Asynchronous Tasks & No & No & Yes & \textbf{Yes} \\ \hline
Native API Security & Yes & Manual & Yes & \textbf{Yes} \\ \hline
\end{tabular}
\end{table}

\section{Research Gaps in Published IEEE Literature}
We used 5 key research streams in the IEEE Xplore online library pertaining to multi-cloud management. The following gaps were found:

\subsection{Gap 1: The "Log Normalization" Deficit}
There is a considerable body of work (e.g., [22,28]) on Multi-Cloud Resource Dis- covery. Nonetheless, the method of normalization of error logs across clouds is not standardized. 

\subsection{Gap 2: Developer Experience (DevEx) in IaC}
Recent studies on IaC (e.g., [5, 26]) are based on the "Language" level. Research on how to make IaC accessible to non-experts by using high-level UI abstractions and background execution is lacking.

\subsection{Gap 3: Automated Secret Injection Security}
In the majority of articles on multi-cloud security (e.g., [27]) a hardware security module (HSM) is assumed. The research work by Nebula investigates on the application-layer symmetric encryption (AES-128) as a tool to securely maintain credential management on portable devices without the need to rely on costly hardware.


\newpage
